\documentclass[11pt]{article}

\usepackage[margin=1in]{geometry}
\usepackage[T1]{fontenc}
\usepackage[utf8]{inputenc}
\usepackage{lmodern}
\usepackage{microtype}
\usepackage{amsmath,amssymb,amsthm}
\usepackage{enumitem}
\usepackage{hyperref}
\usepackage[nameinlink,noabbrev]{cleveref}
\usepackage{booktabs}

\hypersetup{
  colorlinks=true,
  linkcolor=blue,
  citecolor=blue,
  urlcolor=blue
}

\title{Pure Math Benchmark: Toward a Community Standard for Research-Style Mathematical Evaluation}
\author{Xuntao Hu}
\date{March 2026}

\begin{document}
\maketitle

\begin{abstract}
We propose \emph{Pure Math Benchmark}, a public benchmark initiative for evaluating research-style pure mathematical reasoning in language models. Existing public math evaluations are valuable, but they are dominated by short-answer school mathematics, contest-style problem solving, undergraduate coursework coverage, or formal theorem-proving interfaces. None of these paradigms directly addresses broad, natural-language reasoning across modern pure mathematics as practiced by advanced graduate students, PhD researchers, postdocs, and faculty. We therefore propose a benchmark organized by an MSC2020-inspired hierarchy, built around expert-reviewed items, argument-centered scoring, structured provenance metadata, and explicit governance. The benchmark is designed not only as a dataset, but as a program: a public standard with correction, deprecation, editorial review, and community participation pathways. We release an initial seed corpus, a public taxonomy, a per-item schema, and a formal expert-review rubric. We invite mathematicians across subject areas to participate in review, critique, and scope-setting so that this benchmark can mature into a credible shared evaluation resource for mathematical language models.
\end{abstract}

\section{Introduction}

Language models now perform strongly on many public mathematics benchmarks. At the same time, the benchmark ecology for mathematics remains skewed toward tasks that are not adequate proxies for the practice of pure mathematics at research level. A system that performs well on grade-school arithmetic or olympiad-style final-answer tasks need not demonstrate the abilities that matter most to working mathematicians: selecting the right theorem, identifying a missing hypothesis, comparing two frameworks, constructing a counterexample, outlining a viable proof strategy, or making meaningful partial progress on a difficult mathematical question.

This project begins from the claim that \emph{research-style pure mathematics} deserves its own evaluation standard. The aim is not to replace existing benchmarks, but to complement them with a benchmark whose unit of evaluation is the mathematical argument rather than only the final answer.

\section{Why Existing Public Eval Sets Are Not Enough}

Public math evaluation already contains several successful benchmark families, including school-level mathematical QA, contest and olympiad benchmarks, undergraduate subject benchmarks, formal theorem-proving benchmarks, and a small number of newer proof-oriented public evaluations \cite{cobbe2021gsm8k,hendrycks2021measuring,he2024olympiadbench,wang2024omnimath,zheng2025ugmathbench,polu2021formalmath}.

These resources are valuable, but each leaves important gaps for advanced pure mathematics.

\subsection{Final-answer bias}

Many current benchmarks reward the correct endpoint more than the validity of the reasoning. In pure mathematics, this is often the wrong emphasis. A mathematically incorrect proof that lands on the right conclusion should not be treated as equivalent to a valid argument.

\subsection{Contest bias}

Competition mathematics is important and difficult, but it favors polished bounded tasks and trick-sensitive solutions. Research mathematics often involves theorem selection, reformulation, sensitivity to hypotheses, and long-horizon structural reasoning rather than the contest format.

\subsection{Formal-interface mismatch}

Formal theorem-proving benchmarks are indispensable for one capability regime, but they measure proof construction inside formal systems. That is not the same as evaluating natural-language mathematical reasoning as mathematicians currently practice it.

\subsection{Subject imbalance}

Even advanced public math benchmarks frequently overrepresent the areas that are easiest to package as public benchmark problems. Modern pure mathematics is much broader.

\section{Scope and Positioning}

\emph{Pure Math Benchmark} focuses on:
\begin{itemize}[leftmargin=*]
  \item advanced graduate to expert-level pure mathematics,
  \item natural-language reasoning and mathematical exposition,
  \item broad subject coverage using an MSC2020-aware taxonomy,
  \item expert-reviewed items with explicit provenance and contamination metadata.
\end{itemize}

The benchmark does not aim to be:
\begin{itemize}[leftmargin=*]
  \item a generic STEM benchmark,
  \item a contest-only benchmark,
  \item an open-conjecture-solving benchmark,
  \item a formal theorem-proving benchmark in disguise.
\end{itemize}

\section{Taxonomy}

The benchmark uses a three-level hierarchy:
\begin{enumerate}[leftmargin=*]
  \item \texttt{subject\_family},
  \item \texttt{subfield},
  \item \texttt{finer\_topic}.
\end{enumerate}

This public taxonomy is inspired by MSC2020 but made more readable for benchmark users. Each item is also mapped back to one or more official MSC codes \cite{msc2020}.

The initial subject families are:
\begin{itemize}[leftmargin=*]
  \item Logic and Foundations,
  \item Combinatorics and Discrete Structures,
  \item General and Structural Algebra,
  \item Number Theory and Arithmetic Geometry,
  \item Commutative Algebra and Homological Methods,
  \item Representation Theory and Noncommutative Algebra,
  \item Algebraic Geometry,
  \item Topology,
  \item Differential and Symplectic Geometry,
  \item Real, Complex, and Harmonic Analysis,
  \item Functional Analysis and Operator Theory,
  \item Dynamical Systems and Mathematical Physics Interface.
\end{itemize}

\section{Item Design}

The benchmark includes research-style task types such as theorem applicability, proof strategy, missing-hypothesis diagnosis, construction and counterexample, framework comparison, invariant interpretation, and local-to-global reasoning.

Each item carries structured metadata for taxonomy, provenance, review state, leakage assessment, and expected answer format. The repository includes machine-readable item and review schemas to support auditable release discipline.

\section{Expert Review Rubric}

We treat expert review as part of the benchmark definition, not merely as a post hoc annotation layer.

The review rubric tracks:
\begin{itemize}[leftmargin=*]
  \item main claim correctness,
  \item logical validity,
  \item use of hypotheses and definitions,
  \item completeness,
  \item proof strategy quality,
  \item partial progress value,
  \item communication clarity.
\end{itemize}

Responses with fatal mathematical flaws should not be rewarded simply because they contain plausible terminology. Conversely, incomplete but mathematically meaningful progress should be distinguished from bluffing or filler.

\section{Governance}

If a benchmark is meant to become a standard, then governance must be part of the release architecture. We therefore include a Steering Committee model, an Editorial Board model, conflict-of-interest rules, correction and appeals procedures, and versioning rules for benchmark releases.

This framework is designed to prevent the project from becoming an opaque personal dataset.

\section{Seed Release}

The present repository includes:
\begin{itemize}[leftmargin=*]
  \item a public benchmark charter,
  \item comparison documents for the current benchmark landscape,
  \item a benchmark taxonomy and MSC mapping policy,
  \item item and review schemas,
  \item a seed cross-subject corpus,
  \item expert-review rubrics,
  \item governance and participation documents.
\end{itemize}

The seed release is intentionally modest. It is designed to establish standards before scale.

\section{Invitation to the Mathematical Community}

This project will only become useful as a standard if mathematicians help shape it. We invite participation from advanced PhD students, postdoctoral researchers, faculty, subject-matter experts willing to review items, and researchers interested in rigorous evaluation design.

We especially welcome criticism about subject coverage, item quality, benchmark scope, rubric clarity, and governance design.

\section{Feedback and Revision Loop}

The project adopts a public revision loop:
\begin{enumerate}[leftmargin=*]
  \item release benchmark materials publicly,
  \item gather mathematical feedback through issues, discussions, and editorial channels,
  \item route substantive concerns through review and adjudication,
  \item publish corrections and release notes,
  \item revise the manuscript and benchmark artifacts accordingly.
\end{enumerate}

We view this loop as essential. A community standard in mathematics should improve through visible, mathematically serious critique.

\section{Conclusion}

\emph{Pure Math Benchmark} is an attempt to create evaluation infrastructure that better matches the structure of pure mathematics as a research discipline. The benchmark is intentionally broader than contest math, more natural-language oriented than formal theorem proving, and more governance-aware than a one-off dataset release. We hope it can become a useful public instrument for evaluating mathematical language models, and we invite mathematicians to help determine whether it succeeds.

\bibliographystyle{plain}
\bibliography{references}

\end{document}
